\documentclass[10pt,a4paper]{amsart}
\usepackage{amsmath,amssymb,amsthm,graphicx}
% Retained preamble for the Stacks Project chapter snapshot in this directory.
% The chapter's own preamble.tex is a copy of the shared Stacks Project preamble; the block
% below is sliced MECHANICALLY, line for line, from it. Nothing was added, removed or
% redefined; the wrapper class, the xr-hyper/hyperref wiring and the externaldocument list
% are omitted because the assembler supplies the document class.
%
% lines taken from preamble.tex: 3, 5, 9, 11, 12, 156, 157, 158, ... (51 lines total)
%
% package coverage restored from the chapter's own preamble (2026-10-13): xy, xy-2cell

\usepackage[all]{xy}
\xyoption{2cell}
\UseAllTwocells

\newenvironment{reference}{\comment}{\endcomment}
\newenvironment{slogan}{\comment}{\endcomment}
\newenvironment{history}{\comment}{\endcomment}
\theoremstyle{plain}
\newtheorem{theorem}[subsection]{Theorem}
\newtheorem{proposition}[subsection]{Proposition}
\newtheorem{lemma}[subsection]{Lemma}
\theoremstyle{definition}
\newtheorem{definition}[subsection]{Definition}
\newtheorem{example}[subsection]{Example}
\newtheorem{exercise}[subsection]{Exercise}
\newtheorem{situation}[subsection]{Situation}
\theoremstyle{remark}
\newtheorem{remark}[subsection]{Remark}
\newtheorem{remarks}[subsection]{Remarks}
\numberwithin{equation}{subsection}
\def\lim{\mathop{\mathrm{lim}}\nolimits}
\def\colim{\mathop{\mathrm{colim}}\nolimits}
\def\Spec{\mathop{\mathrm{Spec}}}
\def\Hom{\mathop{\mathrm{Hom}}\nolimits}
\def\Ext{\mathop{\mathrm{Ext}}\nolimits}
\def\SheafHom{\mathop{\mathcal{H}\!\mathit{om}}\nolimits}
\def\SheafExt{\mathop{\mathcal{E}\!\mathit{xt}}\nolimits}
\def\Sch{\mathit{Sch}}
\def\Mor{\mathop{\mathrm{Mor}}\nolimits}
\def\Ob{\mathop{\mathrm{Ob}}\nolimits}
\def\Sh{\mathop{\mathit{Sh}}\nolimits}
\def\NL{\mathop{N\!L}\nolimits}
\def\CH{\mathop{\mathrm{CH}}\nolimits}
\def\proetale{{pro\text{-}\acute{e}tale}}
\def\etale{{\acute{e}tale}}
\def\QCoh{\mathit{QCoh}}
\def\Ker{\mathop{\mathrm{Ker}}}
\def\Im{\mathop{\mathrm{Im}}}
\def\Coker{\mathop{\mathrm{Coker}}}
\def\Coim{\mathop{\mathrm{Coim}}}
\DeclareMathSymbol{\boxtimes}{\mathbin}{AMSa}{"02}
\def\QCohstack{\mathcal{QC}\!\mathit{oh}}
\def\Cohstack{\mathcal{C}\!\mathit{oh}}
\def\Spacesstack{\mathcal{S}\!\mathit{paces}}
\def\Quotfunctor{\mathrm{Quot}}
\def\Hilbfunctor{\mathrm{Hilb}}
\def\Curvesstack{\mathcal{C}\!\mathit{urves}}
\def\Polarizedstack{\mathcal{P}\!\mathit{olarized}}
\def\Complexesstack{\mathcal{C}\!\mathit{omplexes}}
\def\Pic{\mathop{\mathrm{Pic}}\nolimits}
\def\Picardstack{\mathcal{P}\!\mathit{ic}}
\def\Picardfunctor{\mathrm{Pic}}
\def\Deformationcategory{\mathcal{D}\!\mathit{ef}}
% lane-s3 patch-sd-envs.tex - DO NOT REMOVE LINES 16-21 (the executable patch).
% Contract (Chief-confirmed 2026-10-14):
%  * Lines 16-18: reference/slogan/history are declared as \comment containers in
%    every retained-preamble of this lane; the patch empties them so navigation
%    pointer text inside statements typesets without "Environment undefined".
%  * Line 21: \cite is globally no-op'd because the corpus rule is "cards carry
%    no bibliography": there is no bibliography backend at assembly time, and
%    every patched card's statement+proof span was scanned zero-cite before
%    patching (2026-10-14, lane-s3) - the only \cite occurrences live inside
%    the emptied reference environments (e.g. \cite{EGA} pointers). Removing
%    line 21 breaks compilation of every card whose statement carries a
%    reference environment with a citation (co01 wave: Citation `EGA' undefined).
%  * The patch is baked into task.json at card-build time (preamble_span
%    end_line must be the last line of this file - the builder computes it
%    dynamically since 2026-10-14); cards without patchenv never load it.
\renewenvironment{reference}{}{}
\renewenvironment{slogan}{}{}
\renewenvironment{history}{}{}
% empty out citations inside navigation-pointer environments (no bibliography
% backend; spans themselves are cite-free - scanned 2026-10-14, lane-s3)
\renewcommand{\cite}[2][]{}
\usepackage{cleveref}
\crefname{originallemma}{original lemma}{original lemmas}
\Crefname{originallemma}{Original lemma}{Original lemmas}
\makeatletter
\newlabel{lemma-compare-torsors}{{04UT}{source}{}{sourcepointer-1}{}}
\newlabel{lemma-compare-torsors@cref}{{[originallemma][0][]04UT}{[1][0][]source}{}{}{}}
\makeatother
\crefname{originallemma}{original lemma}{original lemmas}
\Crefname{originallemma}{Original lemma}{Original lemmas}
\makeatletter
\newlabel{proposition-equal-quotient-stacks}{{04WM}{source}{}{sourcepointer-2}{}}
\newlabel{proposition-equal-quotient-stacks@cref}{{[originallemma][0][]04WM}{[1][0][]source}{}{}{}}
\makeatother
\crefname{originallemma}{original lemma}{original lemmas}
\Crefname{originallemma}{Original lemma}{Original lemmas}
\makeatletter
\newlabel{section-group-quotient-stacks}{{04UV}{source}{}{sourcepointer-3}{}}
\newlabel{section-group-quotient-stacks@cref}{{[originallemma][0][]04UV}{[1][0][]source}{}{}{}}
\makeatother
\crefname{originallemma}{original lemma}{original lemmas}
\Crefname{originallemma}{Original lemma}{Original lemmas}
\makeatletter
\newlabel{spaces-descent-lemma-descending-property-flat}{{041W}{source}{}{sourcepointer-4}{}}
\newlabel{spaces-descent-lemma-descending-property-flat@cref}{{[originallemma][0][]041W}{[1][0][]source}{}{}{}}
\makeatother
\crefname{originallemma}{original lemma}{original lemmas}
\Crefname{originallemma}{Original lemma}{Original lemmas}
\makeatletter
\newlabel{spaces-descent-lemma-descending-property-locally-finite-presentation}{{041T}{source}{}{sourcepointer-5}{}}
\newlabel{spaces-descent-lemma-descending-property-locally-finite-presentation@cref}{{[originallemma][0][]041T}{[1][0][]source}{}{}{}}
\makeatother
\crefname{originallemma}{original lemma}{original lemmas}
\Crefname{originallemma}{Original lemma}{Original lemmas}
\makeatletter
\newlabel{spaces-descent-lemma-descending-property-surjective}{{041Q}{source}{}{sourcepointer-6}{}}
\newlabel{spaces-descent-lemma-descending-property-surjective@cref}{{[originallemma][0][]041Q}{[1][0][]source}{}{}{}}
\makeatother
\crefname{originallemma}{original lemma}{original lemmas}
\Crefname{originallemma}{Original lemma}{Original lemmas}
\makeatletter
\newlabel{spaces-groupoids-definition-principal-homogeneous-space}{{04TY}{source}{}{sourcepointer-7}{}}
\newlabel{spaces-groupoids-definition-principal-homogeneous-space@cref}{{[originallemma][0][]04TY}{[1][0][]source}{}{}{}}
\makeatother
\title{The classifying stack of a group algebraic space}
\author{The Stacks Project Authors}
\date{stacks-project:examples-stacks.tex:local-snapshot-20261012; GNU FDL 1.2 or later}
\begin{document}
\maketitle
\noindent\textit{Editorial note.} Statement span is lines 1675-1693 of sources/stacks-examples-stacks/pinned-original-source.tex; proof span is lines 1695-1725, the complete author proof (29 body lines).\par\medskip
\noindent\textit{Editorial note.} Navigation stubs: 7 cited-but-not-included labels are supplied as non-load-bearing pointers with their pinned tags where the snapshot tag index has them.\par\medskip
\noindent\textit{Editorial note.} extra\_packages is empty on purpose: the retained preamble already loads xy with options and a second declaration clashes in the assembler wrapper.\par\medskip
\noindent\textit{Editorial note.} Text is retained verbatim, including the author's own typos and placeholders.\par\medskip
\noindent\textit{Editorial note.} Queue index 4 of queue\_examples-stacks.json (pb=1463, stmt 1675-1693, proof 1695-1725, halfA).\par\medskip
\noindent\textit{Editorial note.} stacks-examples-stacks wave-1 (exs1): fresh adjudication accept (2026-10-14); pool per Chief 09:4x authorization (idle-line throughput over kun's backup note, SMC precedent); examples-stacks.tex 27 labels/20 with proof judged 9 accept / 11 reject; claim\_key v2 stacks:examples-stacks.tex:<label>. Statement carries a slogan environment (patchenv -> patch-sd-envs.tex).\par\medskip
\noindent\textit{Allowed standard prerequisites (editorial).}\begin{itemize}
\item Explicit prerequisite boundary. The cited inputs are prerequisites; the argument in the retained span is not replaced by any cited result.
\end{itemize}\medskip
\ifdefined\hypertarget\hypertarget{sourcepointer-1}{}\fi
\noindent\textit{Source reference.} Original lemma 04UT: sources/stacks-examples-stacks/pinned-original-source.tex, source lines 1358--1369. Cited by the retained statement or proof but not part of the collected spans; supplied only as a navigation pointer. The argument does not depend on its content.\par\medskip
\ifdefined\hypertarget\hypertarget{sourcepointer-2}{}\fi
\noindent\textit{Source reference.} Original lemma 04WM: sources/stacks-examples-stacks/pinned-original-source.tex, source lines 1528--1536. Cited by the retained statement or proof but not part of the collected spans; supplied only as a navigation pointer. The argument does not depend on its content.\par\medskip
\ifdefined\hypertarget\hypertarget{sourcepointer-3}{}\fi
\noindent\textit{Source reference.} Original lemma 04UV: sources/stacks-examples-stacks/pinned-original-source.tex, source lines 1409--1426. Cited by the retained statement or proof but not part of the collected spans; supplied only as a navigation pointer. The argument does not depend on its content.\par\medskip
\ifdefined\hypertarget\hypertarget{sourcepointer-4}{}\fi
\noindent\textit{Source reference.} Original lemma 041W: sources/stacks-examples-stacks/pinned-original-source.tex, source lines 1725--1725. Cited by the retained statement or proof but not part of the collected spans; supplied only as a navigation pointer. The argument does not depend on its content.\par\medskip
\ifdefined\hypertarget\hypertarget{sourcepointer-5}{}\fi
\noindent\textit{Source reference.} Original lemma 041T: sources/stacks-examples-stacks/pinned-original-source.tex, source lines 1725--1725. Cited by the retained statement or proof but not part of the collected spans; supplied only as a navigation pointer. The argument does not depend on its content.\par\medskip
\ifdefined\hypertarget\hypertarget{sourcepointer-6}{}\fi
\noindent\textit{Source reference.} Original lemma 041Q: sources/stacks-examples-stacks/pinned-original-source.tex, source lines 1725--1725. Cited by the retained statement or proof but not part of the collected spans; supplied only as a navigation pointer. The argument does not depend on its content.\par\medskip
\ifdefined\hypertarget\hypertarget{sourcepointer-7}{}\fi
\noindent\textit{Source reference.} Original lemma 04TY: sources/stacks-examples-stacks/pinned-original-source.tex, source lines 1725--1725. Cited by the retained statement or proof but not part of the collected spans; supplied only as a navigation pointer. The argument does not depend on its content.\par\medskip
\begin{lemma}
\label{lemma-classifying-stacks}
\begin{slogan}
The classifying stack of a group scheme or group algebraic space.
\end{slogan}
Let $S$ be a scheme. Let $B$ be an algebraic space over $S$.
Let $G$ be a group algebraic space over $B$. Then the stacks
in groupoids
$$
[B/G],\quad
[[B/G]],\quad
G\textit{-Torsors},\quad
\mathcal{G}/\mathcal{B}\textit{-Torsors}
$$
are all canonically equivalent.
If $G \to B$ is flat and locally
of finite presentation, then these are also equivalent to
$G\textit{-Principal}$.
\end{lemma}

\begin{proof}
The equivalence
$G\textit{-Torsors} \to \mathcal{G}/\mathcal{B}\textit{-Torsors}$
is given in Lemma \ref{lemma-compare-torsors}.
The equivalence $[B/G] \to [[B/G]]$ is given in
Proposition \ref{proposition-equal-quotient-stacks}.
Unwinding the definition of $[[B/G]]$ given in
Section \ref{section-group-quotient-stacks}
we see that $[[B//G]] = G\textit{-Torsors}$.

\medskip\noindent
Finally, assume $G \to B$ is flat and locally of finite presentation.
To show that the natural functor
$G\textit{-Torsors} \to G\textit{-Principal}$ is an equivalence
it suffices to show that for a scheme $U$ over $B$
a principal homogeneous $G_U$-space $X \to U$
is fppf locally trivial. By our definition of principal homogeneous spaces
(Groupoids in Spaces,
Definition \ref{spaces-groupoids-definition-principal-homogeneous-space})
there exists an fpqc covering $\{U_i \to U\}$ such that
$U_i \times_U X \cong G \times_B U_i$ as algebraic spaces over $U_i$.
This implies that $X \to U$ is surjective, flat, and locally of finite
presentation, see
Descent on Spaces, Lemmas
\ref{spaces-descent-lemma-descending-property-surjective},
\ref{spaces-descent-lemma-descending-property-flat}, and
\ref{spaces-descent-lemma-descending-property-locally-finite-presentation}.
Choose a scheme $W$ and a surjective \'etale morphism $W \to X$.
Then it follows from what we just said that $\{W \to U\}$ is an fppf covering
such that $X_W \to W$ has a section. Hence $X$ is an fppf $G_U$-torsor.
\end{proof}

\end{document}
